
% BEGIN S2
 
\par\medskip\noindent\textbf{2 Preliminaries}\par
 
% BEGIN S2.p1
 
We use the standard notation of equational unification [7] and term rewriting systems [6]. Given a first-order signature $\Sigma$ and a (countable) set of variables $V$, the $\Sigma$-terms over variables $V$ are built in the usual way by taking into account the arity of each function symbol in $\Sigma$. Each $\Sigma$-term is well-formed: if it is rooted by a $n$-ary function symbol in $\Sigma$, then it has necessarily $n$ direct subterms. For any term $t$, $|t|$ denotes the number of symbols occurring in $t$. The set of $\Sigma$-terms over variables $V$ is denoted by $T(\Sigma,V)$. The set of variables from $V$ occurring in a term $t\in T(\Sigma,V)$ is denoted by $\mathit{Var}(t)$. A term $t$ is ground if $\mathit{Var}(t)=\emptyset$. A term $t$ is linear if each variable in $\mathit{Var}(t)$ occurs only once in $t$. For any position $p$ in a term $t$ (including the root position $\varepsilon$), $t(p)$ is the symbol at position $p$, $t|_{p}$ is the subterm of $t$ at position $p$, and $t[u]_{p}$ is the term $t$ in which $t|_{p}$ is replaced by $u$. A substitution is an endomorphism of $T(\Sigma,V)$ with only finitely many variables not mapped to themselves. A substitution is denoted by $\sigma=\{x_{1}\mapsto t_{1},\dots,x_{m}\mapsto t_{m}\}$, where the domain of $\sigma$ is $Dom(\sigma)=\{x_{1},\dots,x_{m}\}$ and the range of $\sigma$ is $Ran(\sigma)=\{t_{1},\dots,t_{m}\}$. Application of a substitution $\sigma$ to $t$ is written $t\sigma$. A $\Sigma$-equation is a pair of $\Sigma$-terms denoted by $s=^{?}t$ or simply $s=t$ when it is clear from the context that we do not refer to an axiom.\par
 
% END S2.p1
 
% BEGIN S2.p2
 
Let $V_{C}$ denote a finite set of context variables (sometimes called holes) such that $V_{C}\cap V=\emptyset$. We use the notation $\diamond_{i}$ to represent the context variable $\diamond_{i}\in V_{C}$. A context is a linear term containing only context variables. More formally, a context is a term, $C\in T(\Sigma,V_{C})$, where each variable (context hole) occurs at most once. Thus, the size of a context follows from the size of a term, where any hole occurrence counts for $1$. Given a context $C$ with $m$ variables and $m$ terms $S_{1},\dots,S_{m}\in T(\Sigma,V)$, $C[S_{1},\dots,S_{m}]$ denotes the term $C\{\diamond_{1}\mapsto S_{1},\dots,\diamond_{m}\mapsto S_{m}\}$ where for $i=1,\dots,m$, $\diamond_{i}$ denotes the $i$-th context variable occurring in the term $C$, and $C$ is called the context part of $C[S_{1},\dots,S_{m}]$. When $S_{1},\dots,S_{m}\in\mathcal{S}$, $C[S_{1},\dots,S_{m}]$ is said to be a context instantiated with terms in $\mathcal{S}$.\par
 
% END S2.p2
 
% BEGIN S2.SS1
 
\par\medskip\noindent\textbf{2.1 Equational Theories}\par
 
% BEGIN S2.SS1.p1
 
Given a set $E$ of $\Sigma$-axioms (i.e., pairs of terms in $T(\Sigma,V)$, denoted by $l=r$), the equational theory $=_{E}$ is the congruence closure of $E$ under the law of substitutivity (by a slight abuse of terminology, $E$ is often called an equational theory). Equivalently, $=_{E}$ can be defined as the reflexive transitive closure $\leftrightarrow_{E}^{*}$ of an equational step $\leftrightarrow_{E}$ defined as follows: $s\leftrightarrow_{E}t$ if there exist a position $p$ of $s$, $l=r$ (or $r=l$) in $E$, and substitution $\sigma$ such that $s|_{p}=l\sigma$ and $t=s[r\sigma]_{p}$. An axiom $l=r$ is regular if $\mathit{Var}(l)=\mathit{Var}(r)$. An axiom $l=r$ is collapse-free if $l$ and $r$ are non-variable terms. An equational theory is regular (resp., collapse-free) if all its axioms are regular (resp., collapse-free). An equational theory $E$ is said to be permutative if for any $l=r$ in $E$, the number of occurrences of any (function or variable) symbol in $l$ is equal to the number of occurrences of that symbol in $r$. Well-known theories such as Associativity ($A=\{(x+y)+z=x+(y+z)\}$), Commutativity ($C=\{x+y=y+x\}$), and Associativity-Commutativity ($AC=A\cup C$) are permutative theories. A theory $E$ is said to be shallow if variables can only occur at a depth at most $1$ in axioms of $E$. For example, $C$ is shallow but $A$ and $AC$ are not. Any constant in $\Sigma$ not occurring in $E$ is said to be free.\par
 
% END S2.SS1.p1
 
% BEGIN Thmtheorem1
 
\par\medskip\noindent\textbf{Definition 1 (Unification).}\par
 
% BEGIN Thmtheorem1.p1
 
Given a signature $\Sigma$, an equational theory $E$, and a set of $\Sigma$-equations, $\Gamma=\{s_{1}=^{?}t_{1},\dots,s_{n}=^{?}t_{n}\}$. The $E$-unification decision problem asks if there exists a substitution $\sigma$ such that $s_{i}\sigma=_{E}t_{i}\sigma$ for all $1\leq i\leq n$. If $E$ is presented as a rewrite relation $R$, then we ask if there exists a substitution $\sigma$ such that $s_{i}\sigma\leftrightarrow_{R}^{*}t_{i}\sigma$ for all $1\leq i\leq n$.\par
 
% END Thmtheorem1.p1
 
% END Thmtheorem1
 
% BEGIN S2.SS1.p2
 
Another important algorithmic problem is matching, which can be seen as a restricted form of the unification problem.\par
 
% END S2.SS1.p2
 
% BEGIN Thmtheorem2
 
\par\medskip\noindent\textbf{Definition 2 (Matching).}\par
 
% BEGIN Thmtheorem2.p1
 
Given a signature $\Sigma$, an equational theory $E$, and a set of $\Sigma$-equations, $\Gamma=\{s_{1}=^{?}t_{1},\dots,s_{n}=^{?}t_{n}\}$. The $E$-matching decision problem asks if there exists a substitution $\sigma$ such that $s_{i}\sigma=_{E}t_{i}$ for all $1\leq i\leq n$. If $E$ is presented as a rewrite relation $R$, then we ask if there exists a substitution $\sigma$ such that $s_{i}\sigma\leftrightarrow_{R}^{*}t_{i}$ for all $1\leq i\leq n$.\par
 
% END Thmtheorem2.p1
 
% END Thmtheorem2
 
% END S2.SS1
 
% BEGIN S2.SS2
 
\par\medskip\noindent\textbf{2.2 Rewrite Relations}\par
 
% BEGIN S2.SS2.p1
 
Given a signature $\Sigma$, an oriented $\Sigma$-axiom is called a rewrite rule of the form $l\rightarrow r$ such that $l,r\in T(\Sigma,V)$, $l$ is not a variable and $\mathit{Var}(r)\subseteq\mathit{Var}(l)$. A finite set of rewrite rules is called a term rewriting system (TRS, for short). Let $R$ be any TRS. For any $\Sigma$-terms $s$ and $t$, $s$ $R$-rewrites to $t$, denoted by $s\rightarrow_{R}t$, if there exist a position $p$ of $s$, $l\rightarrow r\in R$, and substitution $\sigma$ such that $s_{|p}=l\sigma$ and $t=s[r\sigma]_{p}$. The term $s$ is said to be $R$-reducible, $s_{|p}$ is called a redex, and in the particular case where $s_{|p}=l\sigma$, $s$ $R$-rewrites to $t$, denoted by $s\rightarrow_{R}t$. A term is an innermost redex if none of its proper subterms is a redex. The symmetric relation $\leftarrow_{R}\cup\rightarrow_{R}$ is denoted by $\longleftrightarrow_{R}$. The rewrite relation $\rightarrow_{R}$ is confluent if $\longleftrightarrow_{R}^{*}$ is included in $\rightarrow_{R}^{*}\circ\leftarrow_{R}^{*}$. The rewrite relation $\rightarrow_{R}$ is convergent if $\rightarrow_{R}$ is both terminating and confluent. When $\rightarrow_{R}$ is convergent, we have that for any terms $t,t^{\prime}$, $t\longleftrightarrow_{R}^{*}t^{\prime}$ iff $t{\downarrow}_{R}=t^{\prime}{\downarrow}_{R}$, where $t{\downarrow}_{R}$ (resp., $t^{\prime}{\downarrow}_{R}$) denotes the unique normal form of $t$ (resp., $t^{\prime}$) w.r.t $\rightarrow_{R}$. A TRS $R$ is said to be optimally reducing if for any $l\rightarrow r$ in $R$ and any substitution $\theta$, the following holds: ($t\theta$ is $R$-irreducible for any strict subterm $t$ of $l$) implies that $r\theta$ is $R$-irreducible. A TRS is linear if for any $l\rightarrow r$ in $R$, both $l$ and $r$ are linear. A convergent TRS $R$ is said to be subterm convergent if for any $l\rightarrow r\in R$, $r$ is either a strict subterm of $l$ or a constant. We refer to [6] for the classical notion of overlap and critical pair. Hence, a trivial critical pair is obtained by an overlap at the root position of a rule with a renaming of the same rule. A TRS with no non-trivial critical pairs is non-overlapping.\par
 
% END S2.SS2.p1
 
% END S2.SS2
 
% BEGIN S2.SS3
 
\par\medskip\noindent\textbf{2.3 Knowledge Problems}\par
 
% BEGIN S2.SS3.p1
 
The applied pi calculus and frames are used to model attacker knowledge [3]. In this model, the set of messages or terms which the attacker knows, and which could have been obtained from observing one or more protocol sessions, are the set of terms in $Ran(\sigma)$ of the frame $\phi=\nu\tilde{n}.\sigma$, where $\sigma$ is a substitution ranging over ground terms. We also need to model cryptographic concepts such as nonces, keys, and publicly known values. We do this by using names, which are essentially free constants. Here also, we need to track the names which the attacker knows, such as public values, and the names which the attacker does not know a priori, such as freshly generated nonces. $\tilde{n}$ consists of a finite set of restricted names; these names represent freshly generated names which remain secret from the attacker. The set of names occurring in a term $t$ is denoted by $\mathit{fn}(t)$. For any frame $\phi=\nu\tilde{n}.\sigma$, let $\mathit{fn}(\phi)$ be the set of names $\mathit{fn}(\sigma)\backslash\tilde{n}$ where $\mathit{fn}(\sigma)=\bigcup_{t\in Ran(\sigma)}\mathit{fn}(t)$; and for any term $t$, let $t\phi$ denote – by a slight abuse of notation – the term $t\sigma$. For any term $t$, we say that $t$ satisfies the name restriction of $\phi$ if $\mathit{fn}(t)\cap\tilde{n}=\emptyset$.\par
 
% END S2.SS3.p1
 
% BEGIN Thmtheorem3
 
\par\medskip\noindent\textbf{Definition 3 (Deduction).}\par
 
% BEGIN Thmtheorem3.p1
 
Let $\phi=\nu\tilde{n}.\sigma$ be a frame, and $t$ a ground term. We say that $t$ is deduced from $\phi$ modulo $E$, denoted by $\phi\vdash_{E}t$, if there exists a term $\zeta$ such that $\zeta\sigma=_{E}t$ and $\mathit{fn}(\zeta)\cap\tilde{n}=\emptyset$. The term $\zeta$ is called a recipe of $t$ in $\phi$ modulo $E$.\par
 
% END Thmtheorem3.p1
 
% END Thmtheorem3
 
% BEGIN S2.SS3.p2
 
Another form of knowledge is the ability to tell if two frames are statically equivalent modulo $E$.\par
 
% END S2.SS3.p2
 
% BEGIN Thmtheorem4
 
\par\medskip\noindent\textbf{Definition 4 (Static Equivalence).}\par
 
% BEGIN Thmtheorem4.p1
 
Two terms $s$ and $t$ are equal in a frame $\phi=\nu\tilde{n}.\sigma$ modulo an equational theory $E$, denoted $(s=_{E}t)\phi$, if $s\sigma=_{E}t\sigma$, and ${\tilde{n}}\cap(\mathit{fn}(s)\cup\mathit{fn}(t))=\emptyset$. The set of all equalities $s=t$ such that $(s=_{E}t)\phi$ is denoted by $Eq(\phi)$. Given a set of equalities $Eq$, the fact that $(s=_{E}t)\phi$ for all $s=t\in Eq$ is denoted by $\phi\models Eq$. Two frames $\phi=\nu\tilde{n}.\sigma$ and $\psi=\nu\tilde{n}.\tau$ are statically equivalent modulo $E$, denoted as $\phi\approx_{E}\psi$, if $Dom(\sigma)=Dom(\tau)$, $\phi\models Eq(\psi)$ and $\psi\models Eq(\phi)$.\par
 
% END Thmtheorem4.p1
 
% END Thmtheorem4
 
% BEGIN S2.SS3.p3
 
Deduction and static equivalence are known to be decidable for subterm convergent rewrite systems [2].\par
 
% END S2.SS3.p3
 
% END S2.SS3
 
% BEGIN S2.SS4
 
\par\medskip\noindent\textbf{2.4 Graph Embedded TRS}\par
 
% BEGIN S2.SS4.p1
 
It will be useful in one of the dichotomy results that follows to use a restricted class of term rewrite systems for which it is already known that the static equivalence problem is decidable. These systems were introduced in [25] and then simplified in [10], where they are used to expand on the class of subterm convergent term rewrite systems. Actually, we don’t need the full class here for the proofs that follow. Therefore, we only introduce the needed components below, which keeps the definition simpler and more compact. In particular, we introduce a subset of contracting TRSs below that omit the permutative component. One can find the full definition in [10].\par
 
% END S2.SS4.p1
 
% BEGIN S2.SS4.p2
 
The class of contracting TRS first starts by modeling the graph minor relation in the TRS setting. This is done by a set of rewrite schemas. This set of rewrite schemata then induces a graph-embedded term rewrite system. Notice that this is very similar to what is often done when considering the homeomorphic embedding relation in TRSs.\par
 
% END S2.SS4.p2
 
% BEGIN Thmtheorem5
 
\par\medskip\noindent\textbf{Definition 5 (Graph Embedding).}\par
 
% BEGIN Thmtheorem5.p1
 
Consider the following reduction relation, $\rightarrow^{*}_{R_{gemb}}$, where $R_{gemb}$ is the set of rules given by the instantiation of the following rule schema:\par
    \[\left\{\begin{array}[]{ll}&\text{for any }f\in\Sigma\\
(1)&f(x_{1},\ldots,x_{n})\rightarrow x_{i}\\
(2)&f(x_{1},\ldots,\leavevmode\nobreak\ x_{i-1},\leavevmode\nobreak\ x_{i},%
\leavevmode\nobreak\ x_{i+1}\leavevmode\nobreak\ \ldots,\leavevmode\nobreak\ x%
_{n})\rightarrow f(x_{1},\ldots,\leavevmode\nobreak\ x_{i-1},\leavevmode%
\nobreak\ x_{i+1},\leavevmode\nobreak\ \ldots,\leavevmode\nobreak\ x_{n})\\
&\text{and for any }f,g\in\Sigma\\
(3)&f(x_{1},\ldots,x_{i-1},g(\bar{z}),x_{i+1},\leavevmode\nobreak\ \ldots,x_{m%
})\rightarrow g(x_{1},\ldots,\leavevmode\nobreak\ x_{i-1},\bar{z},x_{i+1},%
\leavevmode\nobreak\ \ldots,x_{m})\\
(4)&f(x_{1},\ldots,x_{i-1},g(\bar{z}),x_{i+1},\leavevmode\nobreak\ \ldots,x_{m%
})\rightarrow f(x_{1},\ldots,\leavevmode\nobreak\ x_{i-1},\bar{z},x_{i+1},%
\leavevmode\nobreak\ \ldots,x_{m})\end{array}\right\}\]    
% END Thmtheorem5.p1
 
% BEGIN Thmtheorem5.p2
 
We say a term $t^{\prime}$ is graph-embedded in a term $t$, denoted $t\succcurlyeq_{gemb}t^{\prime}$, if $t^{\prime}$ is a well formed term such that $t\rightarrow_{R_{gemb}}^{*}t^{\prime}$.\par
 
% END Thmtheorem5.p2
 
% BEGIN Thmtheorem5.p3
 
A TRS $R$ is graph-embedded if for any $l\rightarrow r\in R$, $l\succcurlyeq_{gemb}r$ or $r$ is a constant.\par
 
% END Thmtheorem5.p3
 
% END Thmtheorem5
 
% BEGIN S2.SS4.p3
 
The graph embedded class is not sufficient to obtain decidability of static equivalence, so we next need to introduce a type of projection rule that ensures access to subterms and thus decidability of static equivalence.\par
 
% END S2.SS4.p3
 
% BEGIN Thmtheorem6
 
\par\medskip\noindent\textbf{Definition 6 (Projecting Rule).}\par
 
% BEGIN Thmtheorem6.p1
 
Let $R$ be a TRS, $x$ a variable and $t$ a term with a single occurrence of $x$. A projecting rule in $R$ over $t$ leading to $x$ is a rule in $R$ which is a variant of $t^{\prime}\rightarrow x$ where $t^{\prime}$ is a superterm of $t$ with no additional occurrences of $x$.\par
 
% END Thmtheorem6.p1
 
% END Thmtheorem6
 
% BEGIN S2.SS4.p4
 
How these projecting rules relate to the entire term-rewrite system can now be defined.\par
 
% END S2.SS4.p4
 
% BEGIN Thmtheorem7
 
\par\medskip\noindent\textbf{Definition 7 (Projection-Closed Derivation).}\par
 
% BEGIN Thmtheorem7.p1
 
Let $R$ be a TRS. A non-empty $R_{gemb}$-derivation is said to be projection-closed with respect to $R$ if it has the form\par
    \[s\rightarrow_{R_{gemb}}^{l\rightarrow r,\gamma}s^{\prime}\rightarrow_{R_{gemb}%
}^{*}t\]    
where\par
  
\par\noindent  $\blacksquare$ 
% BEGIN S2.I1.i1.p1
 
$s$ is a linear term and $t$ is a well formed term,\par
 
% END S2.I1.i1.p1
 
 
\par\noindent  $\blacksquare$ 
% BEGIN S2.I1.i2.p1
 
$l\rightarrow r$ is a $R_{gemb}$ rule among (1), (2) and (4),\par
 
% END S2.I1.i2.p1
 
 
\par\noindent  $\blacksquare$ 
% BEGIN S2.I1.i3.p1
 
$\gamma$ is a substitution ranging over variables,\par
 
% END S2.I1.i3.p1
 
 
\par\noindent  $\blacksquare$ 
% BEGIN S2.I1.i4.p1
 
if $l\rightarrow r$ is rule (1), $f(x_{1},\ldots,x_{n})\rightarrow x_{i}$, then for any $x_{i}\gamma$ occurring in $t$ there exists a projecting rule in $R$ over $f(x_{1},\ldots,x_{n})\gamma$ leading to $x_{i}\gamma$,\par
 
% END S2.I1.i4.p1
 
 
\par\noindent  $\blacksquare$ 
% BEGIN S2.I1.i5.p1
 
if $l\rightarrow r$ is rule (4), $f(x_{1},\dots,x_{i-1},g(\bar{z}),x_{i+1},\dots,x_{m})\rightarrow f(x_{1},\dots%
,x_{i-1},\bar{z},x_{i+1},\dots,x_{m})$, then for any $z_{i}\gamma$ occurring in $t$ there exists a projecting rule in $R$ over $g(\bar{z})\gamma$ leading to $z_{i}\gamma$,\par
 
% END S2.I1.i5.p1
 
 
\par\noindent  $\blacksquare$ 
% BEGIN S2.I1.i6.p1
 
$s^{\prime}\rightarrow_{R_{gemb}}^{*}t$ is either projection-closed with respect to $R$ or empty.\par
 
% END S2.I1.i6.p1
 
  
% END Thmtheorem7.p1
 
% END Thmtheorem7
 
% BEGIN S2.SS4.p5
 
These rule can be combined into a subset of the graph-embedded TRS, called contracting, for which the knowledge problems are decidable. However, the contracting definition also allows for a type of restricted permutation. For the results in this paper we don’t need to include the permutation component of the contracting definition and by leaving it out we obtain a strict subset of contracting, introduced below, which is simpler and for which the knowledge problems are decidable.\par
 
% END S2.SS4.p5
 
% BEGIN Thmtheorem8
 
\par\medskip\noindent\textbf{Definition 8 (Permutation-free Contracting TRS).}\par
 
% BEGIN Thmtheorem8.p1
 
A TRS $R$ is said to be permutation-free contracting, $\approx$-free contracting for short, if for any non-subterm rule $l\rightarrow r$ in $R$, there exist a position $p$ of $l$, a substitution $\sigma$ ranging over variables, and a projection-closed derivation $g\rightarrow_{R_{gemb}}^{+}d$ with respect to $R$ such that $l{|_{p}}=g\sigma$ and $r=d\sigma$ is of depth $1$.\par
 
% END Thmtheorem8.p1
 
% END Thmtheorem8
 
% BEGIN Thmtheorem9
 
\par\medskip\noindent\textbf{Lemma 9 ( [25, 10]).}\par
 
% BEGIN Thmtheorem9.p1
 
Deduction and static equivalence are decidable for $\approx$-free contracting TRSs.\par
 
% END Thmtheorem9.p1
 
% END Thmtheorem9
 
% BEGIN S2.SS4.SSS0.P0.SPx1
 
\par\medskip\noindent\textbf{Proof.}\par
 
% BEGIN S2.SS4.SSS0.P0.SPx1.p1
 
Follows from [25] where it is shown that the knowledge problems are decidable for contracting TRS, and the fact that $\approx$-free contracting TRS are a strict subset of contracting. $\hfill\blacktriangleleft$\par
 
% END S2.SS4.SSS0.P0.SPx1.p1
 
% END S2.SS4.SSS0.P0.SPx1
 
% END S2.SS4
 
% END S2


% BEGIN S3
 
\par\medskip\noindent\textbf{3 Dichotomy Results}\par
 
% BEGIN S3.p1
 
It is known that Deduction and Static Equivalence are undecidable in general, see [1, 9]. It was also shown in [1] that static equivalence is not reducible to deduction. Interestingly, it is shown in [1] that deduction can be reduced to static equivalence. However, the proof in [1] requires the addition of a free function symbol. It is unknown if the reduction can be done without the addition. Furthermore, there exist theories for which deduction is decidable but static equivalence is not, see, for example, [9].\par
 
% END S3.p1
 
% BEGIN S3.p2
 
In this section we want to compare the knowledge problems to the well-known unification and matching problems. Actually, this work has already been started, motivated by the close similarities in the problems. Previous work has provided dichotomy results that compare a knowledge problem to either unification or matching, demonstrating a theory for which one is decidable but the other is not. Thus, before beginning the new results, let us quickly review these existing results. After the review, we then complete the picture by developing the final needed results comparing the problems. Considering each of the possible cases:\par
  
\par\noindent  1. 
% BEGIN S3.I1.i1.p1
 
Unification vs deduction and static equivalence.\par
  
\par\noindent  (a) 
% BEGIN S3.I1.i1.I1.i1.p1
 
There exist theories for which unification is undecidable but deduction is decidable. It’s shown in [10] that deduction is decidable in permutative theories but in [21] unification is shown to be undecidable for permutative theories.\par
 
% END S3.I1.i1.I1.i1.p1
 
 
\par\noindent  (b) 
% BEGIN S3.I1.i1.I1.i2.p1
 
There exist theories for which unification is decidable but deduction is undecidable. In [5] a theory is developed for which unification is decidable but deduction is undecidable.\par
 
% END S3.I1.i1.I1.i2.p1
 
 
\par\noindent  (c) 
% BEGIN S3.I1.i1.I1.i3.p1
 
There exist theories for which unification is decidable but static equivalence is undecidable. Shown in this paper.\par
 
% END S3.I1.i1.I1.i3.p1
 
 
\par\noindent  (d) 
% BEGIN S3.I1.i1.I1.i4.p1
 
There exist theories for which unification is undecidable but static equivalence is decidable. Shown in this paper.\par
 
% END S3.I1.i1.I1.i4.p1
 
  
% END S3.I1.i1.p1
 
 
\par\noindent  2. 
% BEGIN S3.I1.i2.p1
 
Matching vs deduction and static equivalence.\par
  
\par\noindent  (a) 
% BEGIN S3.I1.i2.I1.i1.p1
 
There exist theories for which matching is undecidable but deduction is decidable. It is shown in [10] that deduction is decidable in permutative theories but in [21] unification is shown to be undecidable for permutative theories. Actually, the problem used in the reduction is a matching problem, so this also demonstrates a theory for which matching is undecidable.\par
 
% END S3.I1.i2.I1.i1.p1
 
 
\par\noindent  (b) 
% BEGIN S3.I1.i2.I1.i2.p1
 
There exist theories for which matching is decidable but deduction is undecidable. This is due to the unification case (1b) above.\par
 
% END S3.I1.i2.I1.i2.p1
 
 
\par\noindent  (c) 
% BEGIN S3.I1.i2.I1.i3.p1
 
There exist theories for which matching is decidable but static equivalence is undecidable. Due to the unification case (1c) above.\par
 
% END S3.I1.i2.I1.i3.p1
 
 
\par\noindent  (d) 
% BEGIN S3.I1.i2.I1.i4.p1
 
There exist theories for which matching is undecidable but static equivalence is decidable. Shown in this paper.\par
 
% END S3.I1.i2.I1.i4.p1
 
  
% END S3.I1.i2.p1
 
  
% END S3.p2
 
% BEGIN S3.p3
 
We now proceed to consider each of the remaining dichotomy problems in turn.\par
 
% END S3.p3
 
% BEGIN S3.SS1
 
\par\medskip\noindent\textbf{3.1 Remaining Dichotomy Problems}\par
 
% BEGIN S3.SS1.SSSx1
 
\par\medskip\noindent\textbf{Problem: Unification Decidable and Static Equivalence Undecidable}\par
 
% BEGIN S3.SS1.SSSx1.p1
 
For this problem, we are comparing the unification problem to the problem of static equivalence, showing that the problems are not comparable. That is, we show that there exist theories for which unification is decidable and static equivalence is not. In particular, we will be using the theory of optimally reducing, linear, and convergent TRS.\par
 
% END S3.SS1.SSSx1.p1
 
% BEGIN S3.SS1.SSSx1.p2
 
In [5] it is shown that there exists an optimally reducing, linear, and convergent TRS for which the cap problem is undecidable. Since the cap problem is a special case of the deduction problem [10], the result developed in [5] provides the following:\par
 
% END S3.SS1.SSSx1.p2
 
% BEGIN Thmtheorem10
 
\par\medskip\noindent\textbf{Lemma 10 ([5]).}\par
 
% BEGIN Thmtheorem10.p1
 
Deduction is undecidable in general for optimally reducing, linear, non-overlapping and convergent TRS.\par
 
% END Thmtheorem10.p1
 
% END Thmtheorem10
 
% BEGIN S3.SS1.SSSx1.p3
 
The result is obtained by using a reduction from the halting problem for 2-counter Minsky machines [5]. We modify that reduction here to extend the result from the deduction problem to the static equivalence problem.\par
 
% END S3.SS1.SSSx1.p3
 
% BEGIN S3.SS1.SSSx1.p4
 
A 2-counter Minsky machine is a counter machine that is equipped with 2 non-negative counters that can be incremented, decremented and checked for equality with zero. Since these machines are Turing complete, the zero halting problem – whether the machine will halt in an accepting state with both counters at 0 – is undecidable. A configuration of a 2-counter machine can be represented as a triple $(C_{1},q_{l},C_{2})$, where $q_{l}$ is the label of the next instruction to be executed and $C_{i}$, $i=1,2$, are the current counter values. As shown in [5] the operations on such a machine can be modeled by a TRS. We modify their construction a little. For any of the finite many states, let $q_{l}$ be a unary function symbol and represent state $l$. Let $f$ and $s$ be unary function symbols, and let $c$ be a rank 3 function symbol. Let $m$ be a nonce and $0$ be a constant representing zero. $s$ is used to represent the successor function, while $c$ is used to encode configurations. For any 2-counter machine, we can encode the initial configuration as $c(s^{k}(0),q_{0}(m),s^{p}(0))$ and the final configuration as $c(0,q_{F}(m),0)$. We also assume that the machine makes no moves once it is in the final state $F$. Then the state transitions of the counter machine can be modeled by the following convergent, linear, and optimally reducing TRS $R$:\par
    $\displaystyle f(c(x,q_{l}(z),y))$ $\displaystyle\rightarrow c(s(x),q_{l+1}(z),y)$     $\displaystyle f(c(x,q_{l}(z),y))$ $\displaystyle\rightarrow c(x,q_{l+1}(z),s(y))$     $\displaystyle f(c(s(x),q_{l}(z),y))$ $\displaystyle\rightarrow c(x,q_{l+1}(z),y)$     $\displaystyle f(c(0,q_{l}(z),y))$ $\displaystyle\rightarrow c(0,q_{l^{\prime}}(z),y))$     $\displaystyle f(c(x,q_{l}(z),s(y)))$ $\displaystyle\rightarrow c(x,q_{l+1}(z),y))$     $\displaystyle f(c(x,q_{l}(z),0))$ $\displaystyle\rightarrow c(x,q_{l^{\prime}}(z),0))$    
% END S3.SS1.SSSx1.p4
 
% BEGIN S3.SS1.SSSx1.p5
 
The $f$ symbol in the above TRS is used to ensure the optimally reducing requirement. In addition, notice that once the final state $F$ is reached the rewrite derivation comes to an end, and this exactly corresponds to the termination of the counter machine. Let $CT$ stand for “configuration terms,” i.e., ground terms of the form $c(s^{i}(0),q,s^{j}(0))$ where $q$ is a state symbol and $S\leavevmode\nobreak\ =\leavevmode\nobreak\ c(s^{k}(0),q_{0}(m),s^{p}(0))$, $T\leavevmode\nobreak\ =\leavevmode\nobreak\ c(0,q_{F}(m),0)$ respectively be the initial and final configuration terms. It can be shown that $T$ can be deduced from $S$ modulo $R$ if and only if $\exists n:f^{n}(S)\;\rightarrow_{R}^{!}\;T$ if and only if the Minsky machine reaches the final configuration.\par
 
% END S3.SS1.SSSx1.p5
 
% BEGIN S3.SS1.SSSx1.p6
 
We now show that the static equivalence problem is undecidable for this theory.\par
 
% END S3.SS1.SSSx1.p6
 
% BEGIN Thmtheorem11
 
\par\medskip\noindent\textbf{Lemma 11.}\par
 
% BEGIN Thmtheorem11.p1
 
The static equivalence problem is undecidable in general for optimally reducing, linear, non-overlapping and convergent TRSs.\par
 
% END Thmtheorem11.p1
 
% END Thmtheorem11
 
% BEGIN S3.SS1.SSSx1.P0.SPx1
 
\par\medskip\noindent\textbf{Proof.}\par
 
% BEGIN S3.SS1.SSSx1.P0.SPx1.p1
 
Let $R$ be the the optimally reducing, linear, and convergent TRS outlined above. Construct two frames as follows:\par
  
\par\noindent  1. 
% BEGIN S3.I2.i1.p1
 
$\phi_{1}=\nu\tilde{n}.\sigma_{1}$, where $\sigma_{1}=\{x_{1}\mapsto c(s^{k}(0),q_{0}(m),s^{p}(0)),\;x_{2}\mapsto T\}$\par
 
% END S3.I2.i1.p1
 
 
\par\noindent  2. 
% BEGIN S3.I2.i2.p1
 
$\phi_{2}=\nu\tilde{n}.\sigma_{2}$, where $\sigma_{2}=\{x_{1}\mapsto c(s^{k}(0),q_{0}(m),s^{p}(0)),\;x_{2}\mapsto d\}$\par
 
% END S3.I2.i2.p1
 
  
where $d$ is a new private name (nonce) in $\tilde{n}$. Thus $\tilde{n}\,=\,\{m,d\}$.\par
 
% END S3.SS1.SSSx1.P0.SPx1.p1
 
% BEGIN S3.SS1.SSSx1.P0.SPx1.p2
 
Notice that the encoding of the initial configuration of the Minsky machine, that the variable $x_{1}$ maps to, is the same in both frames. The only difference between the frames is in the mappings on $x_{2}$. Clearly if $f^{n}(c(s^{k}(0),q_{0}(m),s^{p}(0)))\;\rightarrow_{R}^{!}\;T$ then $x_{2}\;\approx^{?}\;f^{n}(x_{1})$ is unified by $\sigma_{1}$, but not by $\sigma_{2}$.\par
 
% END S3.SS1.SSSx1.P0.SPx1.p2
 
% BEGIN S3.SS1.SSSx1.P0.SPx1.p3
 
The proof in the other direction – that if there is an equation that is unified by $\sigma_{1}$ but not by $\sigma_{2}$ then $T$ can be deduced from $c(s^{k}(0),q_{0}(m),s^{p}(0))$ – can be proved in a way similar to the proof of Proposition 4 (pages 8–9) of [2].\par
 
% END S3.SS1.SSSx1.P0.SPx1.p3
 
% BEGIN ThmClaim1
 
\par\medskip\noindent\textbf{$\blacktriangleright$ Claim 1.}\par
 
% BEGIN ThmClaim1.p1
 
Let $t$ be a ground term and $T$ be a subterm of $t$ at position $p$. If $t$ is reducible, then $t[b]_{p}$ is reducible where $b$ is a free constant.\par
 
% END ThmClaim1.p1
 
% END ThmClaim1
 
% END S3.SS1.SSSx1.P0.SPx1
 
% BEGIN S3.SS1.SSSx1.P0.SPx2
 
\par\medskip\noindent\textbf{Proof.}\par
 
% BEGIN S3.SS1.SSSx1.P0.SPx2.p1
 
Without loss of generality suppose $t\;=\;l\theta$ for left-hand side of $R$. Since there are no moves after the machine reaches the final state, the symbol $q_{F}$ does not appear on any left-hand side. Thus no nonvariable subterm of $t$ is unifiable with $T$. $\hfill\vartriangleleft$ It is also worth noting here that the term rewriting system $R\;\cup\;\{T\,\rightarrow\,b\}$, where $b$ is a free constant, is convergent.\par
 
% END S3.SS1.SSSx1.P0.SPx2.p1
 
% BEGIN S3.SS1.SSSx1.P0.SPx2.p2
 
We can make use of the constant abstraction (“cutting”) function $\mathbf{cut}_{U,d}$ as defined in [2]\par
    $\displaystyle\mathbf{cut}_{T,d}(a)$ $\displaystyle=a\leavevmode\nobreak\ \text{ if $a$ is a name or constant}$     $\displaystyle\mathbf{cut}_{T,d}(g(T_{1},\ldots,T_{k}))$ $\displaystyle=\begin{cases}d\leavevmode\nobreak\ \leavevmode\nobreak\ \text{if%
 }g=c,\leavevmode\nobreak\ c(T_{1},T_{2},T_{3})=_{R}T\\
g(\mathbf{cut}_{T,d}(T_{1}),\ldots,\mathbf{cut}_{T,d}(T_{k}))&\text{otherwise %
}\end{cases}$    
% END S3.SS1.SSSx1.P0.SPx2.p2
 
% BEGIN S3.SS1.SSSx1.P0.SPx2.p3
 
By extension, we also define for any substitution $\sigma$, the substitution $\mathbf{cut}_{T,d}(\sigma)=\{x\mapsto\mathbf{cut}_{T,d}(x\sigma)\}_{x\in Dom(%
\sigma)}$. With $\mathbf{cut}_{T,d}$, we can prove the following:\par
 
% END S3.SS1.SSSx1.P0.SPx2.p3
 
% BEGIN ThmClaim2
 
\par\medskip\noindent\textbf{$\blacktriangleright$ Claim 2.}\par
 
% BEGIN ThmClaim2.p1
 
$M\;\downarrow_{R}\;N$ if and only if $\,\mathbf{cut}_{T,d}(M)\;\downarrow_{R}\;\mathbf{cut}_{T,d}(N)$.\par
 
% END ThmClaim2.p1
 
% END ThmClaim2
 
% END S3.SS1.SSSx1.P0.SPx2
 
% BEGIN S3.SS1.SSSx1.P0.SPx3
 
\par\medskip\noindent\textbf{Proof.}\par
 
% BEGIN S3.SS1.SSSx1.P0.SPx3.p1
 
The “if” part is straightforward. For the “only if” part, suppose\par
    \[\mathbf{cut}_{T,d}(M)\;\not\downarrow_{R}\;\mathbf{cut}_{T,d}(N)\leavevmode%
\nobreak\ \text{ and }\leavevmode\nobreak\ [T/d](\mathbf{cut}_{T,d}(M)\;%
\downarrow_{R}\;[T/d](\mathbf{cut}_{T,d}(N).\]    
Let $M^{\prime}\;=\;\mathbf{cut}_{T,d}(M)\!\downarrow_{R}$ and let $N^{\prime}\;=\;\mathbf{cut}_{T,d}(N)\!\downarrow_{R}$. Now if $T$ is a subterm of $M$ then $d$ occurs in $M^{\prime}$ since every rule in $R$ is variable-preserving. By the previous claim, if $[T/d](M^{\prime})$ is reducible then so is $M^{\prime}$ which leads to a contradiction. $\hfill\vartriangleleft$\par
 
% END S3.SS1.SSSx1.P0.SPx3.p1
 
% BEGIN ThmClaim3
 
\par\medskip\noindent\textbf{$\blacktriangleright$ Claim 3.}\par
 
% BEGIN ThmClaim3.p1
 
Suppose $T$ cannot be deduced from $S=c(s^{k}(0),q_{0}(m),s^{p}(0))$. For any terms $M$ and $N$ such that $(\mathit{fn}(M)\cup\mathit{fn}(N))\cap\tilde{n}=\emptyset$, we have that $(M=_{R}N)\phi_{1}$ implies $(M=_{R}N)\phi_{2}$.\par
 
% END ThmClaim3.p1
 
% END ThmClaim3
 
% END S3.SS1.SSSx1.P0.SPx3
 
% BEGIN S3.SS1.SSSx1.P0.SPx4
 
\par\medskip\noindent\textbf{Proof.}\par
 
% BEGIN S3.SS1.SSSx1.P0.SPx4.p1
 
Assume $M\sigma_{1}\;\downarrow_{R}\;N\sigma_{1}$. By the earlier claim, we have $\mathbf{cut}_{T,d}(M\sigma_{1})\leavevmode\nobreak\ \downarrow_{R}\leavevmode%
\nobreak\ \mathbf{cut}_{T,d}(N\sigma_{1})$. Let us show that $\mathbf{cut}_{T,d}(M\sigma_{1})\;=\;M\,\textbf{cut}_{T,d}(\sigma_{1})$ and $\mathbf{cut}_{T,d}(N\sigma_{1})\;=\;N\,\textbf{cut}_{T,d}(\sigma_{1})$. Suppose by contradiction, say for $M$, that $\mathbf{cut}_{T,d}(M\sigma_{1})\;\neq\;M\,\textbf{cut}_{T,d}(\sigma_{1})$. Then there must be a nonvariable subterm $M^{\prime}$ of $M$ such that $M^{\prime}\sigma_{1}\;\downarrow_{R}\;T$. But this means that $T$ can be deduced from $S$, which is impossible by assumption. Thus, we have $\mathbf{cut}_{T,d}(M\sigma_{1})\;=\;M\,\textbf{cut}_{T,d}(\sigma_{1})$ and $\mathbf{cut}_{T,d}(N\sigma_{1})\;=\;N\,\textbf{cut}_{T,d}(\sigma_{1})$. So, $M\,\textbf{cut}_{T,d}(\sigma_{1})\;\downarrow_{R}\;N\,\textbf{cut}_{T,d}(%
\sigma_{1})$, equivalently $M\sigma_{2}\;\downarrow_{R}\;N\sigma_{2}$. $\hfill\vartriangleleft$\par
 
% END S3.SS1.SSSx1.P0.SPx4.p1
 
% BEGIN S3.SS1.SSSx1.P0.SPx4.p2
 
Thus $\phi_{1}$ and $\phi_{2}$ are statically equivalent if and only if $T$ cannot be deduced from $S$. $\hfill\blacktriangleleft$\par
 
% END S3.SS1.SSSx1.P0.SPx4.p2
 
% BEGIN S3.SS1.SSSx1.P0.SPx4.p3
 
We need now the second component of the result, that unification is decidable for this TRS. We can rely on the following known result:\par
 
% END S3.SS1.SSSx1.P0.SPx4.p3
 
% BEGIN Thmtheorem12
 
\par\medskip\noindent\textbf{Lemma 12 ([22]).}\par
 
% BEGIN Thmtheorem12.p1
 
Unification is decidable for optimally reducing and convergent TRSs.\par
 
% END Thmtheorem12.p1
 
% END Thmtheorem12
 
% BEGIN S3.SS1.SSSx1.P0.SPx4.p4
 
Thus we obtain the desired dichotomy result for this section.\par
 
% END S3.SS1.SSSx1.P0.SPx4.p4
 
% BEGIN Thmtheorem13
 
\par\medskip\noindent\textbf{Theorem 13.}\par
 
% BEGIN Thmtheorem13.p1
 
There exist theories for which unification is decidable but static equivalence is undecidable.\par
 
% END Thmtheorem13.p1
 
% END Thmtheorem13
 
% END S3.SS1.SSSx1.P0.SPx4
 
% BEGIN S3.SS1.SSSx1.P0.SPx5
 
\par\medskip\noindent\textbf{Proof.}\par
 
% BEGIN S3.SS1.SSSx1.P0.SPx5.p1
 
By Lemma 11 and Lemma 12. $\hfill\blacktriangleleft$\par
 
% END S3.SS1.SSSx1.P0.SPx5.p1
 
% END S3.SS1.SSSx1.P0.SPx5
 
% END S3.SS1.SSSx1
 
% BEGIN S3.SS1.SSSx2
 
\par\medskip\noindent\textbf{Problem: Unification Undecidable and Static Equivalence Decidable}\par
 
% BEGIN S3.SS1.SSSx2.p1
 
For this problem we again compare unification to static equivalence. However, now we show that there are theories for which unification is undecidable but static equivalence is decidable. We follow a similar strategy to that in [4], which shows the undecidability of unification for $\Delta$-strong TRSs. We need to modify the reduction so that it models not $\Delta$-strong theories but $\approx$-free contracting . This modification does not impact the unification problem but since the static equivalence problem is decidable for $\approx$-free contracting TRS, we obtain the desired result.\par
 
% END S3.SS1.SSSx2.p1
 
% BEGIN S3.SS1.SSSx2.p2
 
We will use the modified Post Correspondence Problem (MPCP). Let $\Sigma=\{a,b\}$, and let $P=\{(\alpha_{i},\beta_{i})\leavevmode\nobreak\ |\leavevmode\nobreak\ i=1,2,%
\ldots,n\}\subseteq\Sigma^{+}\times\Sigma^{+}$. Then an instance of the MPCP is defined as follows:\par
  
\par\noindent  $\blacksquare$ 
% BEGIN S3.I3.i1.p1
 
Given a non-empty string $\alpha_{0}\in\Sigma^{+}$.\par
 
% END S3.I3.i1.p1
 
 
\par\noindent  $\blacksquare$ 
% BEGIN S3.I3.i2.p1
 
Does there exist indices $j_{1},j_{2},\ldots,j_{m}\in\{1,2,\ldots,n\}$ such that:\par
    \[\alpha_{0}\alpha_{j_{1}},\alpha_{j_{2}}\ldots\alpha_{j_{m}}=\beta_{j_{1}}\beta%
_{j_{2}}\ldots\beta_{j_{m}}\]    
% END S3.I3.i2.p1
 
  
% END S3.SS1.SSSx2.p2
 
% BEGIN S3.SS1.SSSx2.p3
 
To start we need to convert the MPCP into a rewriting problem by interpreting the symbols in $\Sigma$ of the MPCP as unary function symbols, for each $a\in\Sigma$ we have a function $a^{1}()$. Next, for any string $w\in\Sigma^{*}$ we will convert the concatenation operation into function composition, denoted as $\widetilde{w}$ such that:\par
    \[\widetilde{\lambda}(x)=x,\leavevmode\nobreak\ \widetilde{aw}(x)=a(\widetilde{w%
}(x)),\leavevmode\nobreak\ \widetilde{bw}(x)=b(\widetilde{w}(x)).\]    
% END S3.SS1.SSSx2.p3
 
% BEGIN S3.SS1.SSSx2.p4
 
Now, introduce new function symbols, $f^{3}$, and $n$ distinct unary symbols, $g_{1},\ldots,\leavevmode\nobreak\ g_{n}$.\par
 
% END S3.SS1.SSSx2.p4
 
% BEGIN S3.SS1.SSSx2.p5
 
Let $P$ be an instance of the MPCP, from $P$ construct the following TRS, $R_{P}$, For each pair $(\alpha_{i},\leavevmode\nobreak\ \beta_{i})\in P$:\par
    \[R_{P}=R_{P}\cup\{f(\widetilde{\alpha_{i}}(x),g_{i}(y),\widetilde{\beta_{i}}(z)%
)\rightarrow f(x,y,z)\}\]    
These rules can be seen as checking that the strings correspond to a sequence of MPCP blocks. Notice that since the unique $g_{i}$ symbols prevent critical pairs between the rules of $R_{P}$, the TRS is confluent. However, this current system, while convergent, is not a $\approx$-free contracting TRS. Notice that the projection rules of Definition 8 are missing.\par
 
% END S3.SS1.SSSx2.p5
 
% BEGIN S3.SS1.SSSx2.p6
 
Now we need to add a second set of rules, $R_{C}$, that ensure the TRS is $\approx$-free contracting. The key requirement is that for any of the rules in $R_{P}$ that remove the function symbols $a()$, $b()$, and $g_{i}()$, there must exists another rule in $R$ of the form $l[a(x)]\rightarrow x\in R$. Therefore, we introduce $n+2$ new unary function symbols $h_{i}$ for each $i\in\{1,\ldots,n\}$, and $h_{a}$ and $h_{b}$. Define the TRS $R_{C}$ as follows:\par
    $\displaystyle R_{C}$ $\displaystyle=\{h_{i}(g_{i}(y))\rightarrow y\leavevmode\nobreak\ |\leavevmode%
\nobreak\ f(\widetilde{\alpha_{i}}(x),g_{i}(y),\widetilde{\beta_{i}}(z))%
\rightarrow f(x,y,z)\in R_{P}\}\leavevmode\nobreak\ \bigcup$      $\displaystyle\hphantom{==}\{h_{a}(a(x))\rightarrow x,\leavevmode\nobreak\ h_{b%
}(b(x))\rightarrow x\}$    
% END S3.SS1.SSSx2.p6
 
% BEGIN S3.SS1.SSSx2.p7
 
Again notice that due to the unique symbols $h_{i}$, $h_{a}$, and $h_{b}$, there are no critical pairs between the rules of $R_{C}$.\par
 
% END S3.SS1.SSSx2.p7
 
% BEGIN S3.SS1.SSSx2.p8
 
Finally, form the TRS $R_{PC}$ by combining $R_{P}$ and $R_{C}$: $R_{PC}=R_{P}\cup R_{C}$.\par
 
% END S3.SS1.SSSx2.p8
 
% BEGIN Thmtheorem14
 
\par\medskip\noindent\textbf{Lemma 14.}\par
 
% BEGIN Thmtheorem14.p1
 
The TRS $R_{PC}$ is $\approx$-free contracting convergent.\par
 
% END Thmtheorem14.p1
 
% END Thmtheorem14
 
% BEGIN S3.SS1.SSSx2.P0.SPx1
 
\par\medskip\noindent\textbf{Proof.}\par
 
% BEGIN S3.SS1.SSSx2.P0.SPx1.p1
 
The convergence follows from the absence of critical pairs between the rules. $R_{PC}$ is $\approx$-free contracting due to the additional rules in $R_{C}$. $\hfill\blacktriangleleft$\par
 
% END S3.SS1.SSSx2.P0.SPx1.p1
 
% BEGIN Thmtheorem15
 
\par\medskip\noindent\textbf{Lemma 15.}\par
 
% BEGIN Thmtheorem15.p1
 
Static equivalence is decidable for $R_{PC}$.\par
 
% END Thmtheorem15.p1
 
% END Thmtheorem15
 
% END S3.SS1.SSSx2.P0.SPx1
 
% BEGIN S3.SS1.SSSx2.P0.SPx2
 
\par\medskip\noindent\textbf{Proof.}\par
 
% BEGIN S3.SS1.SSSx2.P0.SPx2.p1
 
From Lemma 14 and Lemma 9. $\hfill\blacktriangleleft$\par
 
% END S3.SS1.SSSx2.P0.SPx2.p1
 
% BEGIN Thmtheorem16
 
\par\medskip\noindent\textbf{Theorem 16.}\par
 
% BEGIN Thmtheorem16.p1
 
There exist theories for which unification is undecidable but static equivalence is decidable.\par
 
% END Thmtheorem16.p1
 
% END Thmtheorem16
 
% END S3.SS1.SSSx2.P0.SPx2
 
% BEGIN S3.SS1.SSSx2.P0.SPx3
 
\par\medskip\noindent\textbf{Proof.}\par
 
% BEGIN S3.SS1.SSSx2.P0.SPx3.p1
 
Let $P,\alpha_{0}$ be an arbitrary instance of the MPCP. Let $R_{PC}$ be the TRS constructed as above. From Lemma 15 the static equivalence problem is decidable.\par
 
% END S3.SS1.SSSx2.P0.SPx3.p1
 
% BEGIN S3.SS1.SSSx2.P0.SPx3.p2
 
Now consider the unification problem\par
    \[f(x,y,\widetilde{\alpha_{0}}(x))=_{R_{PC}}f(c,c,c)\]    
where $c$ is a new constant. Notice that due to the string $\alpha_{0}$, we cannot use the rules contained in $R_{C}$ for reducing the left-hand side of this problem to the right, since there are no $h_{i_{j}}$ symbols contained in $\widetilde{\alpha_{0}}$. Thus, the reduction is done via the rules in $R_{P}$. However, as shown in [4] the unification problem above is decidable iff the corresponding MPCP is. $\hfill\blacktriangleleft$\par
 
% END S3.SS1.SSSx2.P0.SPx3.p2
 
% BEGIN Thmtheorem17
 
\par\medskip\noindent\textbf{$\blacktriangleright$ Remark 17.}\par
 
% BEGIN Thmtheorem17.p1
 
It is interesting to note that there is an interesting parallel to the above result already in the literature. As already cited, in [4] it is shown that the unification problem is undecidable in general for $\Delta$-strong theories. In addition, in [19] a procedure is developed for solving the static equivalence problem for $\Delta$-strong theories. However, there is a difference in the definition of $\Delta$-strong from [4] to [19], for example in the first paper the rules in $\Delta$ are required to be dwindling, while in the latter they are not dwindling. Thus, it’s not immediately clear if these two definitions are equivalent.\par
 
% END Thmtheorem17.p1
 
% END Thmtheorem17
 
% END S3.SS1.SSSx2.P0.SPx3
 
% END S3.SS1.SSSx2
 
% BEGIN S3.SS1.SSSx3
 
\par\medskip\noindent\textbf{Problem: Matching Undecidable and Static Equivalence Decidable}\par
 
% BEGIN S3.SS1.SSSx3.p1
 
Next we consider the final dichotomy results that remain, that of matching compared to static equivalence. First, we can note that since we have shown a case for which unification is decidable but static equivalence is undecidable, Theorem 13, we already have the following result due to the fact that if unification is decidable so is matching.\par
 
% END S3.SS1.SSSx3.p1
 
% BEGIN Thmtheorem18
 
\par\medskip\noindent\textbf{Corollary 18.}\par
 
% BEGIN Thmtheorem18.p1
 
There exist theories for which matching is decidable but static equivalence is undecidable.\par
 
% END Thmtheorem18.p1
 
% END Thmtheorem18
 
% BEGIN S3.SS1.SSSx3.p2
 
Next, notice that in the proof of Theorem 16, the right-hand side of the problem, $f(x,y,\widetilde{\alpha_{0}}(x))=_{R_{PC}}f(c,c,c)$ is ground. Therefore, what we have is actually a matching problem and thus we also obtain the following result via the same proof.\par
 
% END S3.SS1.SSSx3.p2
 
% BEGIN Thmtheorem19
 
\par\medskip\noindent\textbf{Corollary 19.}\par
 
% BEGIN Thmtheorem19.p1
 
There exist theories for which matching is undecidable but static equivalence is decidable.\par
 
% END Thmtheorem19.p1
 
% END Thmtheorem19
 
% END S3.SS1.SSSx3
 
% END S3.SS1
 
% END S3


% BEGIN S4
 
\par\medskip\noindent\textbf{4 Decidable Theories}\par
 
% BEGIN S4.p1
 
With the completion of the dichotomy results we next want to consider the question of non-trivial classes of theories for which all the above problems, knowledge, matching, and unification, decidable. One such example is the class of subterm convergent TRSs. These TRSs are very useful in the formal analysis of cryptographic protocols since they can model a number of cryptographic primitives. However, they require that the theory can be oriented into a restricted type of TRS. We introduce in this section another non-trivial example, the separate variable-permuting class of theories which are a restricted form of equational theory. We show that all of the above problems are decidable for this particular class of permutative theories.\par
 
% END S4.p1
 
% BEGIN S4.SS1
 
\par\medskip\noindent\textbf{4.1 Separate Variable-Permuting Theories}\par
 
% BEGIN S4.SS1.p1
 
Let us begin by introducing the class of variable-permuting theories and then the class of separate variable-permuting (SVP) theories.\par
 
% END S4.SS1.p1
 
% BEGIN Thmtheorem20
 
\par\medskip\noindent\textbf{Definition 20 (Variable-Permuting Theory).}\par
 
% BEGIN Thmtheorem20.p1
 
An axiom $l=r$ is said to be variable-permuting [21] if all the following conditions are satisfied:\par
  
\par\noindent  1. 
% BEGIN S4.I1.i1.p1
 
the set of positions of $l$ is identical to the set of positions of $r$,\par
 
% END S4.I1.i1.p1
 
 
\par\noindent  2. 
% BEGIN S4.I1.i2.p1
 
for any non-variable position $p$ of $l$, $l(p)=r(p)$,\par
 
% END S4.I1.i2.p1
 
 
\par\noindent  3. 
% BEGIN S4.I1.i3.p1
 
for any $x\in\mathit{Var}(l)\cup\mathit{Var}(r)$, the number of positions of $x$ in $l$ is identical to the number of positions of $x$ in $r$.\par
 
% END S4.I1.i3.p1
 
  
An equational theory $E$ is said to be separate variable-permuting ($\mathit{SVP}$, for short) if for any $l=r\in E$ we have:\par
  
\par\noindent  $\blacksquare$ 
% BEGIN S4.I2.i1.p1
 
$l=r$ is a variable-permuting axiom,\par
 
% END S4.I2.i1.p1
 
 
\par\noindent  $\blacksquare$ 
% BEGIN S4.I2.i2.p1
 
$l$ and $r$ are rooted by the same function symbol that does not occur elsewhere in $E$.\par
 
% END S4.I2.i2.p1
 
  
An inner constructor symbol of $E$ is a function symbol from $E$ that never occurs at the root of any equality in $E$.\par
 
% END Thmtheorem20.p1
 
% END Thmtheorem20
 
% BEGIN Thmtheorem21
 
\par\medskip\noindent\textbf{Example 21 (Intruder Theory with a Permutative Axiom).}\par
 
% BEGIN Thmtheorem21.p1
 
Let us consider a theory used in practice to model a group messaging protocol [13]. For this protocol, the theory modeling the intruder can be defined [23] as a combination $R\cup K$ where\par
    \[K=\{\mathit{keyexch}(x,\mathit{pk}(x^{\prime}),y,\mathit{pk}(y^{\prime}))=%
\mathit{keyexch}(x^{\prime},\mathit{pk}(x),y^{\prime},\mathit{pk}(y))\}\]    
and\par
    \[R=\left\{\begin{array}[]{ll}\begin{array}[]{lcl}\mathit{adec}(\mathit{aenc}(m,%
\mathit{pk}(sk)),sk)&\rightarrow&m\\
\mathit{checksign}(\mathit{sign}(m,sk),m,\mathit{pk}(sk))&\rightarrow&ok\end{%
array}&\begin{array}[]{lcl}\mathit{getmsg}(\mathit{sign}(m,sk))&\rightarrow&m%
\\
\mathit{sdec}(\mathit{senc}(m,k),k)&\rightarrow&m\end{array}\end{array}\right\}.\]    
% END Thmtheorem21.p1
 
% BEGIN Thmtheorem21.p2
 
$R$ is a subterm convergent TRS and $K$ is a $\mathit{SVP}$ theory sharing with $R$ the constructor symbol $\mathit{pk}$.\par
 
% END Thmtheorem21.p2
 
% END Thmtheorem21
 
% BEGIN S4.SS1.p2
 
The class of $\mathit{SVP}$ theories could be sufficient to specify intruder theories where one needs only a permutation of variables expressed via a single axiom, as in Example 21. Compared to a shallow theory like Commutativity, a $\mathit{SVP}$ theory allows us to express a permutation involving terms with constructor symbols. Thus, $\mathit{SVP}$ theories provide good examples for applying the combination methods developed in [16] to the deduction and static equivalence problems in the union of theories sharing only constructor symbols. We believe that $\mathit{SVP}$ theories could be useful to approximate $AC$, and so to provide an alternative to $AC$ in some cases. The $AC$ theory is clearly very interesting in protocol analysis, for instance to specify the exclusive-or theory, but it is difficult to implement.\par
 
% END S4.SS1.p2
 
% BEGIN S4.SS1.p3
 
Contrary to $AC$, the decision procedures discussed in the following for the deduction and static equivalence problems in $\mathit{SVP}$ theories are easy to implement in existing systems since they are very similar to the ones known in the class of subterm convergent theories, which is ubiquitous in protocol analysis.\par
 
% END S4.SS1.p3
 
% END S4.SS1
 
% BEGIN S4.SS2
 
\par\medskip\noindent\textbf{4.2 Knowledge Problems in Separate Variable-Permuting Theories}\par
 
% BEGIN S4.SS2.p1
 
The deduction problem is known to be decidable for the class of permutative theories [10], and so it is decidable for the class of $\mathit{SVP}$ theories since $\mathit{SVP}$ theories are permutative. The case of static equivalence is more complicated since it is known to be undecidable in permutative theories [15]. Therefore, we need to show that static equivalence becomes decidable when considering the particular case of $\mathit{SVP}$ theories. Static equivalence has been shown to be decidable in shallow permutative theories [17]. However, the $\mathit{SVP}$ theory $K$ is not shallow. As shown below, the approach followed in [17] for the case of shallow permutative theories can be reused for the case of $\mathit{SVP}$ theories. Actually, a decision procedure for static equivalence in any $\mathit{SVP}$ theory $E$ can obtained by computing a finite set of equalities bounded by $|E|$, where $|E|=\max_{\{l\leavevmode\nobreak\ |\leavevmode\nobreak\ l=r\in E\}}|l|$. To define this particular finite set of equalities, we first construct a frame saturation dedicated to $\mathit{SVP}$ theories. Then, the recipes of the (deducible) terms included in that saturation are used to build an appropriate set of bounded equalities.\par
 
% END S4.SS2.p1
 
% BEGIN S4.SS2.p2
 
Let $St(t)$ be the set of subterms of a term $t$ and let $St(\sigma)$ be the set of subterms of a substitution $\sigma$ defined by $St(\sigma)=\bigcup_{x\in Dom(\sigma)}St(x\sigma)$.\par
 
% END S4.SS2.p2
 
% BEGIN Thmtheorem22
 
\par\medskip\noindent\textbf{Definition 22 (Frame Saturation for $\mathit{SVP}$ Theories).}\par
 
% BEGIN Thmtheorem22.p1
 
Let $E$ be any $\mathit{SVP}$ theory. For any frame $\phi=\nu\tilde{n}.\sigma$, let $ESt(\sigma)=St(\sigma)\cup\bigcup_{t\in ET}St(t)$ where $ET$ is the set of terms $s[\vec{d}]$ such that $|s|\leq|E|$, $fn(s)\cap\tilde{n}=\emptyset$, $\vec{d}\in St(\sigma)$ and $s[\vec{d}]\leftrightarrow_{E}^{\epsilon}t$ for some term $t$.\par
 
% END Thmtheorem22.p1
 
% BEGIN Thmtheorem22.p2
 
The set of terms $sat_{E}(\phi)$ is the smallest set $D$ such that:\par
  
\par\noindent  (1) 
% BEGIN S4.I3.i1.p1
 
$Ran(\sigma)\subseteq D$,\par
 
% END S4.I3.i1.p1
 
 
\par\noindent  (2) 
% BEGIN S4.I3.i2.p1
 
if $t_{1},\dots,t_{n}\in D$ and $f(t_{1},\dots,t_{n})\in ESt(\sigma)$ then $f(t_{1},\dots,t_{n})\in D$,\par
 
% END S4.I3.i2.p1
 
 
\par\noindent  (3) 
% BEGIN S4.I3.i3.p1
 
if $t\in D$, $t^{\prime}\in ESt(\sigma)$, $t=_{E}t^{\prime}$, then $t^{\prime}\in D$,\par
 
% END S4.I3.i3.p1
 
  
% END Thmtheorem22.p2
 
% BEGIN Thmtheorem22.p3
 
For any frame $\phi=\nu\tilde{n}.\sigma$, $\sigma_{*}$ is the substitution $\sigma\{\chi_{u}\mapsto u\leavevmode\nobreak\ |\leavevmode\nobreak\ u\in sat_{%
E}(\phi)\backslash Ran(\sigma)\}$ where $\chi_{u}$ is a fresh variable for any $u\in sat_{E}(\phi)\backslash Ran(\sigma)$, and $\phi_{*}$ denotes the frame $\nu\tilde{n}.\sigma_{*}$. Given a recipe $\zeta_{u}$ for each $u\in sat_{E}(\phi)\backslash Ran(\sigma)$, the substitution $\{\chi_{u}\mapsto\zeta_{u}\leavevmode\nobreak\ |\leavevmode\nobreak\ u\in sat_%
{E}(\phi)\backslash Ran(\sigma)\}$ is called a recipe substitution of $\phi$ and is denoted by $\zeta_{\phi}$. Let $Bt=\{t\leavevmode\nobreak\ |\leavevmode\nobreak\ |t|\leq|E|\}$. The set $Eq^{B}(\phi)$ is the set of equalities $t\zeta_{\phi}=t^{\prime}\zeta_{\phi}$ such that $(t\zeta_{\phi}=_{E}t^{\prime}\zeta_{\phi})\phi$ and $t,t^{\prime}\in Bt$.\par
 
% END Thmtheorem22.p3
 
% END Thmtheorem22
 
% BEGIN S4.SS2.p3
 
The two following technical lemmas are similar to the ones developed in [17] for the case of shallow permutative theories.\par
 
% END S4.SS2.p3
 
% BEGIN Thmtheorem23
 
\par\medskip\noindent\textbf{Lemma 23.}\par
 
% BEGIN Thmtheorem23.p1
 
Let $E$ be any $\mathit{SVP}$ theory. For any terms $s$ and $t$ satisfying the name restriction, if $s\phi_{*}=_{E}t\phi_{*}$ and $\psi\models Eq^{B}(\phi)$ then $(s\zeta_{\phi})\psi=_{E}(t\zeta_{\phi})\psi$.\par
 
% END Thmtheorem23.p1
 
% END Thmtheorem23
 
% BEGIN S4.SS2.SSS0.P0.SPx1
 
\par\medskip\noindent\textbf{Proof.}\par
 
% BEGIN S4.SS2.SSS0.P0.SPx1.p1
 
By induction on $\max(|s|,|t|,|E|)-|E|$.\par
  
\par\noindent  $\blacksquare$ 
% BEGIN S4.I4.i1.p1
 
Base case: if $\max(|s|,|t|,|E|)=|E|$, then it is true by definition of $Eq^{B}$.\par
 
% END S4.I4.i1.p1
 
 
\par\noindent  $\blacksquare$ 
% BEGIN S4.I4.i2.p1
 
Inductive step: consider $s$ and $t$ are terms such that $\max(|s|,|t|,|E|)>|E|$, $s=f(s_{1},\dots,s_{n})$, $t=f(t_{1},\dots,t_{n})$, and $s\phi_{*}=_{E}t\phi_{*}$. Then, there exist $f(l_{1},\dots,l_{n})=f(r_{1},\dots,r_{n})\in E$ and substitutions $\varphi,\mu$ such that\par
  
\par\noindent  – 
% BEGIN S4.I4.i2.I1.i1.p1
 
for $i=1,\dots,n$, $s_{i}\phi_{*}=_{E}l_{i}\varphi$ and $t_{i}\phi_{*}=_{E}r_{i}\varphi$. In all these matching problems, note that $l_{i}$ and $r_{i}$ are only built over inner constructors. Thus, these matching problems correspond to syntactic matching. For any variable $x$ in $\mathit{Var}(l_{i})=\mathit{Var}(r_{i})$, only two cases can occur.\par
 
% END S4.I4.i2.I1.i1.p1
 
% BEGIN S4.I4.i2.I1.i1.p2
  
\par\noindent  1. 
% BEGIN S4.I4.i2.I0.i1.I1.i1.p1
 
$x\varphi$ is a term $u[w_{1},\dots,w_{m}]$ where $w_{j}\in Ran(\phi_{*})$ for $j=1,\dots,m$. In that case, we define $x\mu=u[v_{1},\dots,v_{m}]$ where $v_{j}\phi_{*}=w_{j}$ for $j=1,\dots,m$.\par
 
% END S4.I4.i2.I0.i1.I1.i1.p1
 
 
\par\noindent  2. 
% BEGIN S4.I4.i2.I0.i1.I1.i2.p1
 
$x\varphi$ is a subterm $u|_{p}$ ($p\neq\epsilon$) of a term $u\in Ran(\phi_{*})$ such that any subterm $u|_{q}$ for $\epsilon\leq q<p$ is rooted by a inner constructor. Thus, $u|_{p}\in Ran(\phi_{*})$ and so there exists a variable $y$ such that $y\phi_{*}=u|_{p}$. In that case, we define $x\mu=y$.\par
 
% END S4.I4.i2.I0.i1.I1.i2.p1
 
  
With the above definition of $\mu$, we have $\varphi=\mu\phi_{*}$.\par
 
% END S4.I4.i2.I1.i1.p2
 
 
\par\noindent  – 
% BEGIN S4.I4.i2.I0.i2.p1
 
Let $i=1,\dots,n$. By definition of $\mu$, all the function symbol occurrences in $l_{i}\mu$ are already in $s$ or in $l_{i}$. Thus, $|l_{i}\mu|\leq\max(|s_{i}|,|l_{i}|)$. Since $|s_{i}|<|s|$ and $|l_{i}|<|E|$, we have $|l_{i}\mu|<\max(|s|,|E|)$. Using the same argument, for $i=1,\dots,n$, $|r_{i}\mu|<\max(|t|,|E|)$.\par
 
% END S4.I4.i2.I0.i2.p1
 
  
Hence the induction hypothesis applies to $s_{i}\phi_{*}=_{E}(l_{i}\mu)\phi_{*}$ and $t_{i}\phi_{*}=_{E}(r_{i}\mu)\phi_{*}$, implying that\par
  
\par\noindent  – 
% BEGIN S4.I4.i2.I1.i1a.p1
 
$(s_{i}\zeta_{\phi})\psi=_{E}(l_{i}\mu\zeta_{\phi})\psi$ for $i=1,\dots,n$,\par
 
% END S4.I4.i2.I1.i1a.p1
 
 
\par\noindent  – 
% BEGIN S4.I4.i2.I1.i2.p1
 
$(t_{i}\zeta_{\phi})\psi=_{E}(r_{i}\mu\zeta_{\phi})\psi$ for $i=1,\dots,n$.\par
 
% END S4.I4.i2.I1.i2.p1
 
  
Thus we have:\par
     $\displaystyle(s\zeta_{\phi})\psi$ $\displaystyle=$ $\displaystyle f((s_{1}\zeta_{\phi})\psi,\dots,(s_{n}\zeta_{\phi})\psi)$      $\displaystyle=_{E}$ $\displaystyle f((l_{1}\mu\zeta_{\phi})\psi,\dots,(l_{n}\mu\zeta_{\phi})\psi)$      $\displaystyle=_{E}$ $\displaystyle f((r_{1}\mu\zeta_{\phi})\psi,\dots,(r_{n}\mu\zeta_{\phi})\psi)$      $\displaystyle=_{E}$ $\displaystyle f((t_{1}\zeta_{\phi})\psi,\dots,(t_{n}\zeta_{\phi})\psi)$      $\displaystyle=$ $\displaystyle(t\zeta_{\phi})\psi$     
which completes the result.\par
 
% END S4.I4.i2.p1
 
  
% END S4.SS2.SSS0.P0.SPx1.p1
 
% BEGIN S4.SS2.SSS0.P0.SPx1.p2
 
$\hfill\blacktriangleleft$\par
 
% END S4.SS2.SSS0.P0.SPx1.p2
 
% BEGIN Thmtheorem24
 
\par\medskip\noindent\textbf{Lemma 24.}\par
 
% BEGIN Thmtheorem24.p1
 
Let $E$ be any $\mathit{SVP}$ theory. For any frames $\phi$ and $\psi$, $\phi\approx_{E}\psi$ iff $\psi\models Eq^{B}(\phi)$ and $\phi\models Eq^{B}(\psi)$.\par
 
% END Thmtheorem24.p1
 
% END Thmtheorem24
 
% END S4.SS2.SSS0.P0.SPx1
 
% BEGIN S4.SS2.SSS0.P0.SPx2
 
\par\medskip\noindent\textbf{Proof.}\par
 
% BEGIN S4.SS2.SSS0.P0.SPx2.p1
 
(If direction) Let $Eq(\phi)$ be the set of all equalities $s\zeta_{\phi}=t\zeta_{\phi}$ such that $(s\zeta_{\phi}=_{E}t\zeta_{\phi})\phi$.\par
 
% END S4.SS2.SSS0.P0.SPx2.p1
 
% BEGIN S4.SS2.SSS0.P0.SPx2.p2
 
Consider any $s\zeta_{\phi}=t\zeta_{\phi}\in Eq(\phi)$. By definition of $\zeta_{\phi}$ and $\phi_{*}$, we have $(s\zeta_{\phi})\phi=_{E}s\phi_{*}$ and $(t\zeta_{\phi})\phi=_{E}t\phi_{*}$. Since $(s\zeta_{\phi})\phi=_{E}(t\zeta_{\phi})\phi$, we obtain $s\phi_{*}=_{E}t\phi_{*}$. Then, by Lemma 23, we get $(s\zeta_{\phi})\psi=_{E}(t\zeta_{\phi})\psi$, which means that $\psi\models Eq(\phi)$. In a symmetric way, we show that $\phi\models Eq(\psi)$. Then, we can conclude since $\phi\approx_{E}\psi\mbox{ iff }\psi\models Eq(\phi)\mbox{ and }\phi\models Eq(\psi)$. $\hfill\blacktriangleleft$\par
 
% END S4.SS2.SSS0.P0.SPx2.p2
 
% BEGIN Thmtheorem25
 
\par\medskip\noindent\textbf{Theorem 25.}\par
 
% BEGIN Thmtheorem25.p1
 
Deduction and static equivalence are decidable in any $\mathit{SVP}$ theory.\par
 
% END Thmtheorem25.p1
 
% END Thmtheorem25
 
% END S4.SS2.SSS0.P0.SPx2
 
% BEGIN S4.SS2.SSS0.P0.SPx3
 
\par\medskip\noindent\textbf{Proof.}\par
 
% BEGIN S4.SS2.SSS0.P0.SPx3.p1
 
Decidability of deduction follows from the decidability of deduction in permutative theories [10], $\mathit{SVP}$ theories being permutative. Static equivalence is decidable thanks to Lemma 24. $\hfill\blacktriangleleft$\par
 
% END S4.SS2.SSS0.P0.SPx3.p1
 
% BEGIN S4.SS2.SSS0.P0.SPx3.p2
 
The combination method developed in [16] can be directly applied to a union $R\cup E$ where $R$ is a convergent TRS with decidable deduction and static equivalence, and $E$ is a $\mathit{SVP}$ theory. In that case, the function symbols shared by $R$ and $E$ must be constructors for $R$ and inner constructors for $E$. Notice that with a shallow permutative theory $E$, constructor symbols can only occur in the ground terms appearing in $E$, and this is clearly a limitation to the application of the combination method. Thus, considering $\mathit{SVP}$ theories instead of shallow ones allows us to get more significant applications.\par
 
% END S4.SS2.SSS0.P0.SPx3.p2
 
% BEGIN Thmtheorem26
 
\par\medskip\noindent\textbf{Corollary 26.}\par
 
% BEGIN Thmtheorem26.p1
 
Let $R$ be any convergent TRS where both deduction and static equivalence are decidable, and $E$ any $\mathit{SVP}$ theory such that the function symbols shared by $R$ and $E$ are constructors for $R$ and inner constructors for $E$. Then, both deduction and static equivalence are decidable in $R\cup E$.\par
 
% END Thmtheorem26.p1
 
% END Thmtheorem26
 
% END S4.SS2.SSS0.P0.SPx3
 
% END S4.SS2
 
% BEGIN S4.SS3
 
\par\medskip\noindent\textbf{4.3 Unification and Matching in Separate Variable-Permuting Theories}\par
 
% BEGIN S4.SS3.p1
 
Finally, let us consider the unification and matching problems in $\mathit{SVP}$ theories. Similar to the static equivalence case, we cannot rely on the decidability of unification in the whole class of permutative theories, since unification have been shown to be undecidable in permutative theories [26] and even in variable permutative theories [15]. Fortunately, $\mathit{SVP}$ theories are examples of theories with the property of being closed by paramodulation, a notion formally introduced in [24, 20, 18]. Actually, the closure by paramodulation follows from the absence of non-variable overlap between two axioms of any $\mathit{SVP}$ theory. Since unification is decidable and finitary in theories closed by paramodulation [24, 20], we have that unification is finitary in $\mathit{SVP}$ theories. Since $\mathit{SVP}$ theories are regular and collapse-free, any unification algorithm in $\mathit{SVP}$ theories leads to a unification algorithm with free symbols, thanks to the combination method known for unification algorithms in regular collapse-free theories [27], the empty (aka free) theory being regular collapse-free. Consequently, unification with free constants, and so matching is also finitary in $\mathit{SVP}$ theories.\par
 
% END S4.SS3.p1
 
% BEGIN Thmtheorem27
 
\par\medskip\noindent\textbf{Theorem 27.}\par
 
% BEGIN Thmtheorem27.p1
 
Unification and matching are decidable in any $\mathit{SVP}$ theory.\par
 
% END Thmtheorem27.p1
 
% END Thmtheorem27
 
% BEGIN S4.SS3.p2
 
Finally, we can put all the above results together to obtain the following.\par
 
% END S4.SS3.p2
 
% BEGIN Thmtheorem28
 
\par\medskip\noindent\textbf{Theorem 28.}\par
 
% BEGIN Thmtheorem28.p1
 
Unification, matching, deduction, and static equivalence are decidable for any $\mathit{SVP}$ theory.\par
 
% END Thmtheorem28.p1
 
% END Thmtheorem28
 
% END S4.SS3
 
% END S4
